\section{中心估计代入回归后的完整梯度}
\label{est:inner}

回归中心是目标定义的一部分。先证明某个 scalar 对 KL 无偏，再把它塞进平方损失，不能自动获得正确回归梯度。固定 anchor $q$，记 $U_a=\log(p_a/q_a)$、$f_a$ 冻结、$z_a=U_a-\eta f_a$，score 为 $\psi_a=\nabla\log p_a$。

\begin{proposition}[独立 reward center 的 Mean 梯度]
\label{est:mean-center}
设随机 $\widehat c_f$ 与主 action $A\sim q$ 独立，分布不依赖 student 参数，二阶矩有限，反传时冻结，$\E\widehat c_f=\E_q[f]$。则
\begin{equation}
\widehat L_M=\frac12[U_A-\eta(f_A-\widehat c_f)]^2
\end{equation}
满足
\begin{equation}
\E\widehat L_M=L_M+\frac{\eta^2}{2}\Var(\widehat c_f),\qquad
\E\nabla\widehat L_M=\nabla L_M.
\end{equation}
\end{proposition}
\begin{proof}
令 $R_M=U_A-\eta(f_A-\E_q[f])$、$\varepsilon=\widehat c_f-\E_q[f]$。展开
\begin{equation}
\widehat L_M=\frac12R_M^2+\eta R_M\varepsilon+\frac{\eta^2}{2}\varepsilon^2.
\end{equation}
独立性和 $\E\varepsilon=0$ 使交叉项期望为零。由于 reward center 的分布与 student 参数无关，最后的方差项没有参数导数。也可直接求导：$\nabla\widehat L_M=(R_M+\eta\varepsilon)\psi_A$，独立零均值误差使额外项期望为零。
\end{proof}

\begin{theorem}[BPO 的独立辅助 center 与成对表示]
\label{est:bpo-sampling}
令 $A,B_1,\ldots,B_m$ 独立同分布于 $q$，$m\geq1$，$\bar z_B=m^{-1}\sum_jz_{B_j}$。则
\begin{equation}
\E[(z_A-\bar z_B)^2]=(1+m^{-1})\Var_q(z).
\end{equation}
因此
\begin{equation}
\widehat L_B=\frac{m}{m+1}\frac{(z_A-\bar z_B)^2}{2\eta}
\end{equation}
对 BPO loss 与其直接 autograd 梯度均无偏。对于两个独立 $A,B\sim q$，等价的成对形式为
\begin{equation}
L_B=\frac1{4\eta}\E[(z_A-z_B)^2],\qquad
\nabla L_B=\frac1{2\eta}\E[(z_A-z_B)(\psi_A-\psi_B)].
\end{equation}
\end{theorem}
\begin{proof}
记 $c_z=\E_q[z]$。独立性给出 $\E(z_A-c_z)(\bar z_B-c_z)=0$，且 $\Var(\bar z_B)=\Var(z_A)/m$。故
\begin{equation}
\E(z_A-\bar z_B)^2=\Var(z_A)+\Var(\bar z_B)=(1+m^{-1})\Var(z_A).
\end{equation}
有限词表上的有限求和可与微分交换，$q$ 冻结，因此修正后的梯度期望同样无偏。$m=1$ 时即得 pairwise loss。直接求导 $\nabla(z_A-z_B)=\psi_A-\psi_B$，乘平方项的导数得到梯度公式。展开成对期望，
\begin{equation}
\E[(z_A-z_B)(\psi_A-\psi_B)]
=2\{\E_q[z\psi]-\E_q[z]\E_q[\psi]\},
\end{equation}
右侧为 $2\eta\nabla L_{\mathrm{BPO}}$，与成对平方损失的 $1/(2\eta)$ 因子相乘恢复 BPO 梯度。
\end{proof}

若同 prefix 有 $m\geq2$ 个 actions，leave-one-out center 使用其余 $m-1$ 个样本，每个主 action 与自己的 center 独立。因此平方项修正系数为 $(m-1)/m$。不同主 action 的 estimator 相互相关，不能把总方差简单除以 $m$。另一种等价无偏目标为成对 U-statistic
\begin{equation}
\widehat L_{B,U}=\frac1{4\eta\binom m2}\sum_{i<j}(z_{A_i}-z_{A_j})^2.
\end{equation}
每一对的期望为 $4\eta L_B$，故平均无偏；固定 $q$ 下反传也无偏。

\subsection{不同信息需求的等价表示}
\label{est:bpo-access}
若 student logits 为 $h_a$，则
\begin{equation}
U_a-U_b=h_a-h_b-\log q_a+\log q_b.
\end{equation}
student log-normalizer 精确抵消；$z_a-z_b$ 也不含 $\log Z_f$、reward center 或 student KL center。因此 pairwise BPO 可以用两个同-prefix teacher 查询和相应 student/anchor 值实现。center 表示则适合完整 logits 或可缓存/可重算的条件积分；两者的查询成本和方差应独立比较。

若 $A,B$ 独立采样于行为分布 $\beta$，且覆盖 $q$ 的支持，令 $w_a=q_a/\beta_a$ 在本轮冻结，则
\begin{equation}
\widehat L_{B,\beta}=\frac{w_Aw_B}{4\eta}(z_A-z_B)^2.
\end{equation}
逐项展开双重有限求和，$\beta_aw_a=q_a$、$\beta_bw_b=q_b$，故期望等于 $L_B$。冻结权重的梯度也无偏。若先把单 token $z_a$ 乘重要性权重再平方，则得到 $w_a^2z_a^2$，不能恢复同一目标。

\subsection{冻结 partition 误差与动态 center 误差}
\label{est:center-error}
设 Part 使用整轮冻结 $\widehat{\log Z_f}=\log Z_f+\delta_Z$。残差变成 $R_P+\delta_Z$，梯度变化为
\begin{equation}
\widehat g_P-g_P=\delta_Z \E_q[\psi].
\end{equation}
anchor 下 $\E_q[\psi]$ 一般不为零。这是改变冻结回归 target 的精确影响；在 $p=q$ 的 inner 起点暂时为零，不能因此宣布整个 inner 都无偏。

若 Mean 的 reward center 改为 $\E_q[f]+\delta_f$，同样有
\begin{equation}
\widehat g_M-g_M=\eta\delta_f \E_q[\psi].
\end{equation}
若中心与主 action 相关，梯度变化为 $\eta\E_q[\delta_f(A)\psi_A]$，不能把误差均值单独提出。

对 BPO，精确 student center 为 $K_S=\KL(q\Vert p)$。若 PG 实现使用每步重算但权重中 detach 的 $K_S+\delta(\theta)$，其他量精确，则
\begin{equation}
\widehat g_B-g_B=\frac{\delta(\theta)}\eta \E_q[\psi]
=-\frac{\delta(\theta)}\eta\nabla K_S.
\end{equation}
其反对称 Jacobian 为
\begin{equation}
\partial_i\widehat g_{B,j}-\partial_j\widehat g_{B,i}
=-\eta^{-1}[(\partial_i\delta)(\partial_jK_S)-(\partial_j\delta)(\partial_iK_S)].
\end{equation}
原 BPO Hessian 对称，$K_S$ 的 Hessian 对称，故只剩误差和 center 梯度的外积。若两梯度不平行，近似 PG 向量场非保守。直接对含近似 center 的平方损失求导则必然保守，却一般得到另一个梯度。两条实现路径在中心近似后应分别验证。
